\documentclass[11pt,a4paper]{article}
\usepackage[margin=27mm]{geometry}
\usepackage{fontspec}
\setmainfont{Latin Modern Roman}
\usepackage{amsmath,amssymb,amsthm,mathtools}
\usepackage{mathrsfs}
\usepackage{longtable,booktabs,array,enumitem,microtype}
\usepackage{xurl}
\usepackage[unicode,colorlinks=true,linkcolor=blue,citecolor=blue,urlcolor=blue]{hyperref}
\usepackage{bookmark}
\newtheorem{theorem}{Statement}
\newtheorem{proposition}[theorem]{Statement}
\newtheorem{lemma}[theorem]{Statement}
\newtheorem{corollary}[theorem]{Statement}
\newtheorem{definition}[theorem]{Statement}
\newtheorem{remark}[theorem]{Statement}
\newtheorem{obstruction}[theorem]{Statement}
\newcommand{\Gal}{\operatorname{Gal}}
\newcommand{\Cent}{\operatorname{Cent}}
\newcommand{\Norm}{\operatorname{Norm}}
\newcommand{\Aut}{\operatorname{Aut}}
\newcommand{\Hom}{\operatorname{Hom}}
\newcommand{\im}{\operatorname{im}}
\newcommand{\Q}{\mathbb Q}
\newcommand{\Z}{\mathbb Z}
\newcommand{\F}{\mathbb F}
\newcommand{\Gm}{\mathbb G_m}
\newcommand{\Zhat}{\widehat{\mathbb Z}}
\newcommand{\Galois}{\Gamma}
\newcommand{\Gzero}{\Gamma^{\circ}}
\newcommand{\Cyc}{K(\mu_\infty)}
\newcommand{\sep}{\mathrm{sep}}
\newcommand{\ab}{\mathrm{ab}}
\newcommand{\tors}{\mathrm{tors}}
\newcommand{\et}{\mathrm{\acute et}}
\newcommand{\status}[1]{\textbf{#1.}}
\newcommand{\passage}[1]{\par\smallskip\noindent\textit{Manuscript: #1}\par\smallskip}
\newcommand{\resref}[1]{\hyperref[res:#1]{#1}}
\newcommand{\Rtwo}{\resref{A6}}
\newcommand{\Arith}{\resref{A1}}
\makeatletter
\renewcommand\paragraph{\@startsection{paragraph}{4}{\z@}%
  {2.7ex \@plus 0.7ex \@minus 0.2ex}{0.5ex \@plus 0.2ex}%
  {\normalfont\normalsize\bfseries}}
\makeatother
\setlength{\emergencystretch}{2em}
\title{The cyclotomic kernel and cuspidal centralizers\\\large Proof audit and calculations for cote 149}
\author{AI-produced by GPT-6 Astra\\OpenAI, ChatGPT Pro mode}
\date{12 November 2026}
\begin{document}
\maketitle
\begin{abstract}
We prove the cyclotomic-kernel centralizer theorem with two cusp-uniqueness arguments, audit the arithmetic and birational comparison maps, and give explicit covering and nilpotent-quotient calculations. Corrections to the manuscript are distinguished from transcription evidence. The coverage ledger records the remaining obligations, including the full local-field calculation and the unrestricted action question. This AI-produced proof audit makes no claim of priority or of a proof of a section conjecture.
\end{abstract}
\input{sections/conventions.tex}
\clearpage
\tableofcontents
\section{Scope and claim ledger}
Statement identifiers A1--E4 and AE1--AE4 are stable identifiers for the accompanying results table. An Open statement gives its proved portion and an explicit remaining obligation. The report distinguishes these from requested work still marked ``unfinished''. None of the claims labelled as unresolved is used as a hypothesis of a proof. The supplied research prompt, prior AI drafts and handwritten scans are evidence to audit, rather than mathematical authorities.
\section{A. Arithmetic and recognition of inertia}
\IfFileExists{sections/A_agent.tex}{\input{sections/A_agent.tex}}{}
\IfFileExists{sections/A_root.tex}{\input{sections/A_root.tex}}{}
\section{B. Arithmetic centralizers}
\IfFileExists{sections/B_agent.tex}{\input{sections/B_agent.tex}}{}
\section{C. Consequences for cote 149}
\IfFileExists{sections/C_root.tex}{\input{sections/C_root.tex}}{}
\IfFileExists{sections/B_birational.tex}{\input{sections/B_birational.tex}}{}
\section{D. Fixed points and cyclotomic base fields}
\IfFileExists{sections/D_root.tex}{\input{sections/D_root.tex}}{}
\IfFileExists{sections/D_agent.tex}{\input{sections/D_agent.tex}}{}
\section{E. Other subgroups and local fields}
\IfFileExists{sections/E_arithmetic_agent.tex}{\input{sections/E_arithmetic_agent.tex}}{}
\IfFileExists{sections/E_root.tex}{\input{sections/E_root.tex}}{}
\IfFileExists{sections/E_local_agent.tex}{\input{sections/E_local_agent.tex}}{}
\IfFileExists{sections/E4_prior_hunches.tex}{\input{sections/E4_prior_hunches.tex}}{}
\nocite{Cote149,LongueMarche,GrothendieckFaltings1983}
\bibliographystyle{alpha}
\bibliography{CYCLOTOMIC_KERNEL_references}
\end{document}
